\documentclass[UTF8,a4paper,11pt]{ctexart}
\usepackage{amsmath,amssymb}
\usepackage[margin=2.5cm]{geometry}
\usepackage[hidelinks]{hyperref}
\title{2010 年北京高考理科第 18 题\\含参对数函数的单调性}
\date{2026 年 10 月 24 日}
\begin{document}
\maketitle
\noindent\textbf{原题来源：}\url{https://github.com/DxAThing/Gaokao-Math-Problems-Compilation/blob/f1ae169b58a75b3db6e35c670b3f0d5c78718188/content/2010/beijing_science.tex#L352}
\par
函数 $f_k(x)=\ln(1+x)-x+\frac{k}{2}x^2$，其中 $k\ge0$，
定义域为 $(-1,+\infty)$。

\paragraph{切线}
\[
 f_k'(x)=\frac1{1+x}-1+kx
        =\frac{x(kx+k-1)}{1+x}.
\]
当 $k=2$ 时，$f_2(1)=\ln2$，$f_2'(1)=\frac32$。
在该点的切线方程为
\[
 \boxed{y-\ln2=\frac32(x-1)}.
\]

\paragraph{完整参数分类}
导数的分母 $1+x$ 恒为正。令 $r=(1-k)/k$（$k>0$）。
\begin{enumerate}
\item 当 $k=0$ 时，$f_k'=-x/(1+x)$。函数在
$(-1,0)$ 严格递增，在 $(0,+\infty)$ 严格递减。
\item 当 $0<k<1$ 时，$r>0$，导数在
$(-1,0),(0,r),(r,+\infty)$ 上依次为正、负、正。
故前、后两个区间递增，中间递减。
\item 当 $k=1$ 时，$f_1'=x^2/(1+x)\ge0$，
且除 $x=0$ 外严格正。因此函数在整个
$(-1,+\infty)$ 上严格递增。
\item 当 $k>1$ 时，$-1<r<0$，导数在
$(-1,r),(r,0),(0,+\infty)$ 上依次为正、负、正。
故前、后两个区间递增，中间递减。
\end{enumerate}

\noindent\textbf{核验边界：}数学推导独立复核；
新源码尚未通过 XeLaTeX 编译验收。
\end{document}
